(a) The evaluation map $\mathcal E\to\mathcal E^{\vee\vee}$ sends $e$ to $(\lambda\mapsto\lambda(e))$. On an open set where $\mathcal E$ is free, it sends a basis to its double dual basis, hence is an isomorphism.

(b) The map $\mathcal E^\vee\otimes\mathcal F\to\mathcal Hom_{\mathcal O_X}(\mathcal E,\mathcal F)$ sends $\lambda\otimes s$ to $(e\mapsto\lambda(e)s)$. If $e_1,\ldots,e_n$ is a local basis, its inverse sends $u$ to $\sum_i e_i^\vee\otimes u(e_i)$.

(c) Send $u:\mathcal E\otimes\mathcal F\to\mathcal G$ to the map $s\mapsto(e\mapsto u(e\otimes s))$. Conversely, send $v:\mathcal F\to\mathcal Hom(\mathcal E,\mathcal G)$ to $e\otimes s\mapsto v(s)(e)$. These constructions commute with restriction and are inverse.

(d) The adjunction map $f^{-1}f_*\mathcal F\to\mathcal F$ induces a natural map $f_*\mathcal F\otimes\mathcal E\to f_*(\mathcal F\otimes f^*\mathcal E)$. Locally on $Y$, write $\mathcal E\cong\mathcal O_Y^n$. This map becomes the natural identification $(f_*\mathcal F)^n\cong f_*(\mathcal F^n)$, so it is an isomorphism.
